\textbf{Schedule-based Tensor Optimizations.}
Schedule-based tensor compilers (e.g., TVM~\cite{tvm} and Halide~\cite{halide})
separate the mathematical definition of a tensor program from \mbox{optimization} decisions.
Optimization decisions are conveyed via dedicated scheduling languages,
with a library of program transformations, e.g., for loops,
parallelization, and data movement.
%A {schedule} is a sequence of transformations applied to the program to optimize it.
%
Auto-schedulers (e.g., Ansor~\cite{ansor} and Halide's auto-sche\-duler~\cite{halide-autosched})
aim to relieve users from hand-writing schedules by searching over the space of
legal transformation sequences.
They rely on heuristic rules and cost models to prune unfavorable schedules.
The design of schedule languages adds to the difficulty of auto-scheduling:
they tend to expose a set of fine-grained transformations
and mix optimizations from different levels (e.g., high-level fusion decisions with
low-level thread mappings) \mbox{in a single schedule.}

\textbf{Tile Optimizations.}
Tile-based compilers (e.g., Triton~\cite{triton} and TileLang~\cite{tilelang})
are an alternative approach to effectively automate low-level optimizations.
Programmers describe the workload as tile computation (fixed-shape subtensors)
in a tile-based language with an SIMD-like programming model,
and the compiler maps the tiles to hardware SIMD units.
While tile compilers automate SIMD-level optimizations,
they leave most higher-level optimizations, e.g., operator fusion and expression rewriting,
to the programmer.

\textbf{Bridging Schedule and Tile Optimizations.}
Prior work (e.g., Neptune~\cite{neptune}) has shown that schedule-based and tile-based
approaches are complementary and can combine to achieve higher performance than either.
Schedule languages are well-suited for high-level optimizations such as operator fusion,
and tile compilers excel at low-level mapping of tile programs to hardware.
Neptune proposes a unified framework where scheduling acts on loop-based programs
(loop-scalar IR in Neptune), which Neptune translates into a tile-based representation (tile IR)
to exploit Triton for SIMD optimizations.
Neptune still requires major manual effort to use, as it does not include an auto-scheduler
and requires the user to provide a schedule.
As Neptune only targets Triton as its optimizer, it cannot make use of
features on newer GPUs such as asynchronous tensor core matmul,
and delivers subpar performance compared to newer kernel libraries.
% \sasa{IMPORTANT: Why Neptune doesn't wor with modern GPUs or backends!!}
% \sasa{Also why it doesn't autoschedule.}
% \matei{However, Neptune is complicated to use since it does not come with an
% auto-scheduler. Additionally, it only targets the Triton backend, which isn't able to
% make use of features on newer GPUs such as asynchronous memory loads, preventing
% performance from being competitive with newer systems.}

\textbf{Tensor Operator Fusion.}
In GPU kernels, multiple outermost loops execute in parallel across different GPU thread blocks.
Combining compute operations under the same parallel loops can bring significant
performance benefit,
since it reduces data communication across thread blocks. This technique is called operator fusion.
Operator fusion is not always possible for computation with complex data dependency patterns.
Neptune generalizes operator fusion to support fusing multiple reduction operations,
and provides them as program transformation passes that \NAME{} uses.

% Additionally, fusing deeper-nested for loops can also be advantageous, as it can
% improve cache locality of data access and allow for better use of shared memory, which
% is a limited resource.
% Being able to fuse these outer loops thus allows for multiple computations to be
% batched together in the same GPU kernel launch, which eliminates the expensive cost of
% launching a kernel.
